\documentclass[a4paper,11pt]{article}
\usepackage[scale={0.8,0.9},centering,includeheadfoot]{geometry}
\usepackage{ifpdf}
\usepackage{amstext}
\usepackage{amsmath}
\usepackage{verbatim}
\newcommand{\vect}[1]{\boldsymbol{#1}}
\begin{document}
\section*{dGompertz}

\subsection*{Parametrisation}

The Gompertz distribution with cure fraction has log ``survial function''
\begin{displaymath}
    \log S(y) = -\frac{\mu}{\alpha}\left(\exp(\alpha y) -1\right)
\end{displaymath}
for response $y\ge 0$, $\mu>0$. The ``cummulative distribution
function'' and the ``density'' then follows as
\begin{displaymath}
    F(y) = 1 - \exp\left[ -\frac{\mu}{\alpha}\left(\exp(\alpha y)
        -1\right) \right]
\end{displaymath}
and
\begin{displaymath}
    f(y) = \mu \exp\left[ \alpha y -\frac{\mu}{\alpha}\left(\exp(\alpha y)
        -1\right) \right].
\end{displaymath}
Note that $\alpha$ is allowed to be negative making to allow for a
cure fraction, hence $\lim_{y\rightarrow\infty} F(y) < 1$ for
$\alpha<0$.

\subsection*{Link-function}
The parameter $\mu$ is linked to the linear predictor $\eta$ as:
\[
    \mu = \exp(\eta)
\]

\subsection*{Hyperparameters}

The shape parameter $\alpha$ is represented as
\begin{displaymath}
    \alpha = \theta
\end{displaymath}
and the prior is defined on $\theta$. 

\subsection*{Specification}

\begin{itemize}
\item \texttt{family="dgompertzsurv"} for survival models
\item Required arguments: $y$ (to be given in a format by using
    $\texttt{inla.surv()}$ for survival models )
\end{itemize}

\subsubsection*{Hyperparameter spesification and default values}
%% DO NOT EDIT!
%% This file is generated automatically from models.R
\begin{description}
	\item[doc] \verb!destructive gompertz (survival) distribution!
	\item[hyper]\ 
	 \begin{description}
	 	\item[theta]\ 
	 	 \begin{description}
	 	 	\item[hyperid] \verb!108101!
	 	 	\item[name] \verb!shape!
	 	 	\item[short.name] \verb!alpha!
	 	 	\item[output.name.intern] \verb!alpha_intern for dGompertz!
	 	 	\item[output.name] \verb!alpha parameter for dGompertz!
	 	 	\item[initial] \verb!-1!
	 	 	\item[fixed] \verb!FALSE!
	 	 	\item[prior] \verb!normal!
	 	 	\item[param] \verb!0 10!
	 	 	\item[to.theta] \verb!function(x) x!
	 	 	\item[from.theta] \verb!function(x) x!
	 	 \end{description}
	 \end{description}
	\item[experimental] \verb!TRUE!
	\item[survival] \verb!TRUE!
	\item[discrete] \verb!FALSE!
	\item[link] \verb!default log neglog!
	\item[pdf] \verb!dgompertz!
\end{description}



\subsection*{Example}

In the following example we estimate the parameters in a simulated
case%%
\verbatiminput{example-dgompertz.R}

\subsection*{Notes}

\end{document}


%%% Local Variables: 
%%% TeX-master: t
%%% End: 
